\documentclass[../../main.tex]{subfiles}

\begin{document}

\section{$L^{1}$ 中的弱序列紧性: Dunford-Pettis 定理}\label{section:ch19-dunford-pettis}

与 $1<p<\infty$ 的情形不同, $L^{1}(X,\mu)$ 一般不是自反的, 因此其闭单位球不是弱紧的. 然而, Dunford-Pettis 定理给出了 $L^{1}$ 中弱序列紧性的刻画: 一致可积性.

\begin{theorem}[Dunford-Pettis]\label{theorem:ch19-dunford-pettis}
令 $(X,\mathcal{M},\mu)$ 为有限测度空间, 而 $\{f_{n}\}$ 是 $L^{1}(X,\mu)$ 中的序列. 则以下陈述等价:
\begin{description}
\item[（i）] $\{f_{n}\}$ 是一致可积的(关于 $\mu$).
\item[（ii）] $\{f_{n}\}$ 的每个子序列有弱收敛的子序列, 即 $\{f_{n}\}$ 是相对弱序列紧的.
\end{description}
\end{theorem}

\begin{proof}
我们概述证明. 令 $\{f_{n}\}\subset L^{1}(X,\mu)$ 为一致可积的. 对每个 $E\in\mathcal{M}$, 定义 $\nu_{n}(E)=\int_{E}f_{n}\,d\mu$. 则每个 $\nu_{n}$ 是关于 $\mu$ 绝对连续的有限带号测度. 一致可积性意味着 $\{\nu_{n}\}$ 关于 $\mu$ 一致绝对连续. 由于 $\{\|f_{n}\|_{1}\}$ 是有界的, $\{\nu_{n}(X)\}$ 也是有界的. 根据 Vitali-Hahn-Saks 定理(定理\cref{theorem:ch18-vitali-hahn-saks}), 存在子序列 $\{\nu_{n_{k}}\}$ 和有限带号测度 $\nu$, 使得对所有 $E\in\mathcal{M}$, $\nu_{n_{k}}(E)\to\nu(E)$. 且 $\nu\ll\mu$.

根据 Radon-Nikodym 定理, 存在 $f\in L^{1}(X,\mu)$ 使得对所有 $E\in\mathcal{M}$, $\nu(E)=\int_{E}f\,d\mu$. 因此对所有可测集 $E$, $\int_{E}f_{n_{k}}\,d\mu\to\int_{E}f\,d\mu$. 由简单函数逼近, 这对所有 $g\in L^{\infty}(X,\mu)$ 成立, 即
\begin{align*}
\int_{X}f_{n_{k}}g\,d\mu\to\int_{X}fg\,d\mu\text{ 对所有 }g\in L^{\infty}(X,\mu).
\end{align*}
因此 $\{f_{n_{k}}\}$ 弱收敛于 $f$.

逆方向: 若 $\{f_{n}\}$ 是相对弱序列紧的, 则它是一致可积的. 这是因为弱收敛序列是有界的, 且若不一致可积, 则存在 $\varepsilon>0$ 和可测集序列 $E_{k}$ 使得 $\mu(E_{k})\to 0$ 但 $\int_{E_{k}}|f_{n_{k}}|\,d\mu\geqslant\varepsilon$. 取 $g_{k}=\operatorname{sgn}(f_{n_{k}})\chi_{E_{k}}\in L^{\infty}$, 则 $\|g_{k}\|_{\infty}\leqslant 1$ 且 $\int f_{n_{k}}g_{k}\,d\mu\geqslant\varepsilon$, 与弱收敛矛盾.
\end{proof}

\begin{remark}
Dunford-Pettis 定理是 $L^{1}$ 中弱紧性的基本刻画. 它告诉我们, 在有限测度空间上, $L^{1}$ 中的相对弱紧集恰是一致可积集. 这与 $L^{p}$($1<p<\infty$) 中简单的有界性不同.
\end{remark}

\begin{exercise}\label{exercise:ch19-11}
令 $(X,\mathcal{M},\mu)$ 为有限测度空间. 证明: $\{f_{n}\}\subset L^{1}(X,\mu)$ 是一致可积的当且仅当它是相对弱紧的.
\end{exercise}

\begin{exercise}\label{exercise:ch19-12}
给出一个 $L^{1}([0,1],m)$ 中的有界序列, 它没有弱收敛的子序列.（提示: 考虑 $f_{n}=n\cdot\chi_{[0,1/n]}$.）
\end{exercise}

\end{document}